\explainedfigure{fig_measured_perception_degradation}{0.97}{Measured effect of simulated fog and rain on the camera perception}{Left: the share of vehicles in view that the YOLO object detector finds, per weather condition (vertical lines are 95\% intervals). Right: the share of vehicle pixels at which the stereo-depth algorithm returns a valid distance. Weather conditions are the teammates' image augmentation at the intensity shown (0 = none, 1 = strongest). 40 scenes; 72 vehicle observations in clear weather and 144 in each weather condition.}{In clear weather the detector finds 72\% of the vehicles. In fog of intensity 1.0 it finds 31\%, in rain of intensity 0.6 it finds 15\% and in rain of intensity 1.0 only 1.4\%. Stereo depth stays valid for 88\% of pixels in clear weather and 49\% in the strongest fog, but remains at 67\% to 83\% in rain.}{The perception measurement confirms that the simulated weather is a real degradation of the camera pipeline, mostly through \emph{missed vehicles}. For vehicles that are detected, the distance estimate is not worse. This tells us what a driving model would have to cope with, and explains why the driving results in the next figures are unaffected.}{fig:perception}
